\begin{abstract}
\noindent
The Erdős--Straus conjecture asserts that $4/n=1/x+1/y+1/z$ has a solution in positive integers for every $n\ge2$. This paper sets out, with proofs, what two bodies of research on the conjecture have established, and adds new results. The first is the Erdős--Straus workbench of The Clankers: a 619-page archive, a supplement of 72 statements on bounded congruence transport, a 124-page reader and a continuation of 28 items. The second is the author's notes of April--June 2026.

The research is organised by the shell criterion, which reduces solvability at a prime $p$ to the occupancy of finitely many shells, each decided by a divisor congruence. On this basis the paper proves exact criteria for the first two shells, a family of residue barriers culminating in the width-one barrier of the notes, and a converse under an exponent condition. It determines the record primes of the least-class function up to $3\cdot10^9$, proves a general fixed-divisor construction of which the archive's four explicit families are instances, and shows that the archive's carrier census can never force the rows that are squares, $1/1024$ of its final domain. From the notes it proves that their terminal core is exactly the set of classes that their certificates cannot reach, and it derives a quadratic-residue sieve on a counterexample whose survivors have density zero.

The research also builds bridges from the equation to other areas, and four of the results of this paper hold in those areas independently of the equation. The notes place the shadow of their order-$7$ mixed mock vector, whose functions agree with Ramanujan's order-$7$ mock theta functions, on three weight-$\tfrac32$ unary theta series; the paper proves that these carry exactly the modular representation of the minimal model $M(2,7)$ after conjugation and a twist by the character of $\eta^3$, and that the analogous order-$5$ block carries the Fibonacci modular data. It classifies the orbits of the workbench's $(3,4,\infty)$ monodromy covers at every level by one invariant. And it shows that the modular flow of the state on $M_2(\C)$ with density proportional to $\operatorname{diag}(e^r,e^{-r})$ and the Lorentz boost of the notes' Lorentzian geometry are the two real forms of one complex one-parameter subgroup of $SL_2(\C)$. Each bridge is described with what is established about it, and the negative results are stated with their exact scope. The workbench does not claim a proof of the conjecture, and nothing here proves or disproves it.
\end{abstract}

\section*{Introduction}\phantomsection\label{sec:about}
\addcontentsline{toc}{section}{Introduction}

\subsection*{The problem and this record}

The conjecture has been verified for all $n\le10^{17}$ \cite{Salez}, it holds for almost all $n$ \cite{Vaughan}, and classical polynomial identities solve it outside the six classes of primes that are squares modulo $840$ \cite{ElsholtzTao,Mordell}. What remains is a question about those classes, and on them polynomial identities provably cannot help (Section~\ref{sec:problem}).

The workbench of The Clankers attacks this question through the shell criterion. Fix a prime $p$ and a candidate least denominator $a$; the rest of the equation is solvable exactly when $a^2$ has a divisor in one of two residue classes modulo $R=4a-p$. The workbench builds on this a large calculus of sufficient conditions, obstructions and residue censuses, and around it geometric and operator-theoretic constructions. The author's notes of April--June 2026 develop the same shell arithmetic in a vocabulary of supports and attach structural layers to it: Niemeier and umbral data, the class $6B$ of the Monster, Ogg's primes, Lorentzian geometry, Cayley--Dickson algebras and mock theta functions.

The workbench and the notes are research carried out with AI systems. The author used them to formulate and test many hypotheses quickly, across several directions at once. Some directions are finished theorems, some are programmes with proved parts, and some are proposals. Not every direction is expected to hold. An unfinished programme is still mathematics, however, and its proved parts stand on their proofs whatever becomes of the rest.
This paper is meant to let a reader see all of this without working through more than a thousand pages. It states what holds, with proofs or the exact location of a proof, what fails and in which scope, and how the research connects with other areas. The audit behind it had limited compute and did not examine every direction. Each direction is described on its own merit and for its bearing on the equation, and where the paper gives a view rather than a proof, the view is labelled as mine and its reasons are given.

\subsection*{Main results}

Throughout, $p$ is an odd prime. For $a>p/4$ with $p\nmid a$ put $R_a=4a-p$ and
\[
E_a=\{u: u\mid a^2,\ R_a\mid 4u+1\},\qquad M_a=\{u: u\mid a^2,\ R_a\mid u+a\};
\]
the shell $a$ is \emph{occupied} if $E_a\cup M_a\ne\emptyset$.

\begin{introthm}[Theorem~\ref{thm:shell}, Propositions~\ref{prop:groupbarrier} and~\ref{prop:groupconverse}]
(i) The equation is solvable at $p$ if and only if some shell $a$ with $p/4<a<p/2$ is occupied, and the ordered pairs $(y,z)$ completing a shell $a$ number exactly $2|E_a|+|M_a|$.
(ii) Let $p\equiv1\pmod4$, $a=(p+R)/4$ with $\gcd(R,pa)=1$, and let $H_a\le(\ZZ/R)^\times$ be the subgroup generated by the primes dividing $a$. If neither $-1$ nor $-4$ lies in $H_a$, the shell $a$ is empty. If $2v_\ell(a)\ge\operatorname{ord}_R(\ell)-1$ for every prime $\ell\mid a$, the shell is occupied if and only if $-1$ or $-4$ lies in $H_a$.
\end{introthm}

Part (i) is the archive's divisor-shell bijection with its layer count, completed here by the range of the least denominator. Part (ii) is the sharpest of a family of residue barriers: the Jacobi-symbol barrier of the archive, a group barrier, and the width-one barrier of the notes. It is a shell form of the quadratic-character mechanism behind the obstruction to polynomial identities, and it says what a proof must produce: for some $R$, a factorisation of $(p+R)/4$ whose prime factors generate $-1$ or $-4$ modulo $R$.

Following Ionascu and Wilson \cite{IonascuWilson}, let $h(p)$ be the least $i$ such that some solution at $p$ has least denominator $x\le(p+4i-1)/4$; a prime is a \emph{strict record} if $h$ is smaller at every smaller prime.

\begin{introthm}[Proposition~\ref{prop:records}]
The strict records of $h$ up to $3\cdot10^9$ are exactly $2$, $73$, $1129$, $1201$, $21{,}169$, $118{,}801$ and $8{,}803{,}369$, with $h=1$, $2$, $3$, $6$, $8$, $15$ and $27$. At $8{,}803{,}369$ the least occupied shell has $R=107$, and exactly one unordered solution has that least denominator.
\end{introthm}

These seven are the records printed by Ionascu and Wilson, recovered in the archive with one class label corrected; that there is no further record up to $3\cdot10^9$ is not in the archive or in their paper, and it was computed by two independent programs. The archive's four explicit families solve the equation on progressions that meet every hard class (Theorem~\ref{thm:families}), and they are instances of a general fixed-divisor construction (Proposition~\ref{prop:fixeddivisor}). Every carrier of the archive's census is of this kind, and none reaches a class that is a square, so the census can never force its square rows, $1/1024$ of its final domain (Proposition~\ref{prop:nosquare}, Corollary~\ref{cor:censuslimit}).

The author's notes contribute a terminal core and a system of certificates. Put $M_{23}=840\cdot11\cdot19\cdot23$, let $S_{840}$ be the six unit squares modulo $840$, and let $Q_q$ be the non-zero squares modulo $q$.

\begin{introthm}[Theorems~\ref{thm:termcore} and~\ref{thm:eightshifts}, Proposition~\ref{prop:densityzero}]
(i) Among the $23{,}760$ classes modulo $M_{23}$ that lie over $S_{840}$ and are units modulo $11$, $19$ and $23$, those on which no elementary certificate of the notes, with residual built from $3,5,7,11,19,23$, applies on any lift are exactly the $2970$ classes of $S_{840}\times Q_{11}\times Q_{19}\times Q_{23}$. The single identities of Salez's seven families leave $3520$ classes open at this level: the core and $550$ further classes.
(ii) A counterexample $p\equiv1\pmod{24}$ is a quadratic residue modulo every odd prime factor of $(p+1)(p+2)(p+3)(p+4)(p+7)(p+8)(2p+1)(3p+1)$.
(iii) The primes $p\equiv1\pmod{24}$ that satisfy the condition from $p+1$ alone have relative density zero among all primes $p\equiv1\pmod{24}$.
\end{introthm}

One half of (i) is proved here and the other is a finite computation, reproduced by an independent program. The conditions from $p+2$, $p+8$ and $2p+1$ are derived here from the notes' certificate branches and are not stated in the archive. At the record prime $8{,}803{,}369$ the other six shifts give no certificate, and those from $p+2$ and $2p+1$ certify it. With certificates of three further shapes, only $23$ primes $p\equiv1\pmod{24}$ below $10^7$ have no certificate of the kinds considered.

The notes attach to the cubic middle of the shell problem an explicit rank-three mixed mock modular object, whose functions agree with Ramanujan's order-$7$ mock theta functions through $q^{40}$ (a reading pass), and they place its shadow on three weight-$\tfrac32$ unary theta series. For $r\in\{1,5,11\}$ let $G_r(\tau)=\sum_{\ell\in\ZZ}v_r(\ell\bmod84)\,\ell\,q^{\ell^2/168}$ be the weight-$\tfrac32$ unary theta series on the notes' three support vectors $v_r$ (Section~\ref{ssec:mock}), let $\rho$ be the pair of matrices by which $(G_1,G_5,G_{11})$ transforms under $S$ and $T$, and let $\rho_{2,7}$ be the representation of $SL_2(\ZZ)$ on the characters of the minimal model $M(2,7)$.

\begin{introthm}[Theorem~\ref{thm:m27}]
There is a signed permutation matrix $P$ with $\overline{\rho(S)}=P\rho_{2,7}(S)P^{-1}$ and $\overline{\rho(T)}\,e(1/8)=P\rho_{2,7}(T)P^{-1}$. So $\bar\rho\otimes\varepsilon^3\cong\rho_{2,7}$, where $\varepsilon$ is the multiplier of Dedekind's $\eta$: the theta series of the notes' shadow, conjugated and twisted by the character of $\eta^3$, transform exactly as the characters of $M(2,7)$.
\end{introthm}

The identity is verified exactly in $\Q(\zeta_{168})$. For order $5$ the same method gives the Fibonacci modular data (Proposition~\ref{prop:fibonacci}). Finally, a package of the workbench builds covers of the projective line from the integral monodromy of the $S^6$ period system; they have a section exactly when the equation is solvable. Write a point of the affine hyperplane $H_D$ of the dual lattice modulo $D$ as $\hat\gamma+b\hat u+c\hat w+d\hat\delta$.

\begin{introthm}[Theorem~\ref{thm:orbits}]
For every level $D\ge1$ the orbits of the monodromy group on $H_D$ are the sets $\{\gcd(b,c,\gcd(D,6))=g\}$, $g\mid\gcd(D,6)$, and they coincide with the orbits of the full stabilizer of the isotropic vector $\gamma$.
\end{introthm}

\subsection*{What the results mean}

My assessment is that, taken together, the results locate the difficulty of the conjecture precisely, for the following reasons. No polynomial identity solves the equation on a residue class that is a square \cite{ElsholtzTao}, and no fixed-divisor certificate reaches such a class (Proposition~\ref{prop:nosquare}). The archive's carrier census runs into this obstruction: its carriers force $97.606\%$ of its final domain and never a square row (Corollary~\ref{cor:censuslimit}). What decides the square classes is the factorisation of the numbers $(p+R)/4$: by Theorem~I(ii) a shell can be occupied only if the prime factors of $(p+R)/4$ generate $-1$ or $-4$ modulo $R$, and under the exponent condition this is also sufficient. A proof of the conjecture therefore has to control, for some $R$, the multiplicative structure of $(p+R)/4$ (Question~\ref{q:global}).

The quadratic-residue sieve shows how restrictive this already is. A counterexample must be a quadratic residue modulo every odd prime factor of eight shifted forms and several more. The primes that satisfy the first condition alone have density zero, and only $23$ primes $p\equiv1\pmod{24}$ below $10^7$ have no certificate of the kinds considered. Whether infinitely many primes have none is Question~\ref{q:qrsurvivors}; a negative answer would prove the conjecture at all but finitely many primes.

The structural results are of a second kind. Theorem~IV places the shadow of the notes' order-$7$ vector inside the modular representation of the $(2,7)$ minimal model. Through the agreement of that vector with Ramanujan's order-$7$ mock theta functions, the first of which gives the $\widehat Z$ invariant of the Brieskorn sphere $-\Sigma(2,3,7)$ \cite{CFS}, the statement concerns these classical objects, and it is independent of the Erdős--Straus problem. A statement about the equation from such structures needs a typed map from the shell data to them; those maps are the subject of Section~\ref{sec:overlaps} and of Questions~\ref{q:modularobject} and~\ref{q:thin}.

\subsection*{The bridges}

In the author's view the connections this research makes between the equation and other areas are among its most valuable outcomes. Section~\ref{sec:overlaps} treats each with the typed map, what is established, what is not verified here, and my assessment. Four results of this paper belong to the other areas themselves and hold independently of the equation: Theorem~IV, Proposition~\ref{prop:fibonacci}, Theorem~V, and Proposition~\ref{prop:flowboost}, by which the modular flow of the state on $M_2(\C)$ with density proportional to $\operatorname{diag}(e^r,e^{-r})$ and the Lorentz boost of the notes' Lorentzian geometry are the compact and the non-compact real forms of one complex one-parameter subgroup of $SL_2(\C)$. A companion note takes up the question this raises for the Apollonian group, the standard thin group acting on the same Minkowski space \cite{ApollonianNote}. Other bridges are partly constructed, such as the split-zero support states and the umbral and Ogg structures of the residue layer, and the bridge to computability through the $\Pi_1$ form of the conjecture is exact. The edges to the author's other workbenches (the zeta programme, $S^6$, Navier--Stokes and Collatz) are listed with their status.

\subsection*{How this record was made}

This record is an audit, carried out on 26--27 September 2026, of the workbench's published documents and of the author's notes of April--June 2026. The workbench is the work of The Clankers, the collective author name of KokunoYumeto and the contributors credited in the workbench's \texttt{README.md} and \texttt{CONTRIBUTIONS.md}:
\begin{itemize}
\item u/CommonCareful3149, for the mod-$107$ arithmetic-progression observation;
\item u/UmbrellaCorp\_HR, for arithmetic, positivity-code and prime-production contributions and for the central-angle and inversive-circle theorems of the archive's §6;
\item u/CivQ17, for the diagrams behind result R05;
\item u/Far-Lifeguard519, for the initiating computational investigation;
\item u/TextBackground496 and u/JGPTech, for source ideas;
\item the owner, KokunoYumeto (u/lepthymo), for the research directions and examples, for directing and developing the investigations with AI, and for the proof and source-fidelity requirements.
\end{itemize}
The archive records language and computational models as contributing-model provenance for particular calculations, not as authors. The notes are the author's.

The audit was carried out by Claude (Anthropic), model \texttt{claude-opus-5-5}, which prepared this paper (\emph{claude-ab} below). Independent Claude instances, each of which wrote none of the text, contributed as follows: a first-pass audit of the workbench; referee passes, which read for overstatement and for understatement and missed generality and found several of the results stated here, each credited at its statement; and six reading passes over the author's notes. Every result found by another instance was verified by claude-ab before it was applied, except the statements marked \emph{(a reading pass)}: those were checked by program in one of the reading passes and not re-derived here, and they count as \emph{not verified here}.

Every statement has one of four statuses:
\begin{itemize}
\item \emph{verified};
\item \emph{not verified here}, which implies nothing in either direction;
\item \emph{my assessment}, a view in the first person with its reasons, introduced by those words;
\item \emph{proved negative}, stated with its counterexample or derivation; the negative results on the workbench's mechanisms are collected in Section~\ref{sec:negative}, and those on the notes are stated where they occur.
\end{itemize}
The status lines under the results make the first two precise. \textsf{Proved here} means the proof is given in full; \textsf{Generalised here} means that the workbench's statement is proved here under weaker hypotheses; \textsf{Audited} means the workbench's proof was read in full, every step checked, and independent checks run; \textsf{Audited (certificate)} means the finite certificate on which a proof rests was recomputed independently and the deduction from it read; \textsf{Reproduced} means the finite content of a statement was recomputed by an independent program and its general proof not re-derived. These are forms of \emph{verified}. \textsf{Audited in part} is verified in the parts named. \textsf{Located} means the statement and its proof were found and are described, and it is a form of \emph{not verified here}. The workbench's own labels (proved, certificate-checked, conjectural, withdrawn) are quoted where they differ.

The workbench was read at commit \texttt{d20b32e} of \texttt{KokunoYumeto/erdos-straus-foundation}, which also holds the 488-page reader of 7 August and its source, and in the Zenodo collection 10.5281/zenodo.20401937, whose record 10.5281/zenodo.22678971 holds the 619-page archive and its source. Locators have the form \emph{archive §$n$, p.~$m$} for the 619-page archive, \emph{supplement §$n$} for the 86-page supplement, and file names and line numbers for the repository; the notes are cited by title and theorem number. Version~3 of this paper and the audit notes cited as \note{NN\_} were withdrawn from the branch \texttt{claude/claude-ab-grind-20260925} on 27 September 2026 because of their framing. They remain, unchanged, in the branch's history at commit \texttt{a90d9e5b}, in the folder \texttt{contrib/claude-ab/grind\_20260925/}, where the folders cited as \note{checks/\ldots} also are. The scripts on which this paper's statements rest are republished in \texttt{contrib/claude-ab/rewrite/checks/es/} on the same branch (Appendix~\ref{app:verification}).

This is version~4. It keeps the theorems, proofs and citations of version~3 and corrects its framing; like the author's latest record, it does not treat the drafts that record supersedes. It adds the results marked as new in version~4: Proposition~\ref{prop:densityzero}, Theorem~\ref{thm:m27}, Propositions~\ref{prop:fibonacci}, \ref{prop:dossier}, \ref{prop:torsor} and~\ref{prop:flowboost}, and Theorem~\ref{thm:orbits} with its corollary. Re-deriving the statements of version~3 from their sources changed some of them. In particular, the identification of the notes' shadow with $M(2,7)$ holds in the form of Theorem~IV, where version~3 said that it fails, and the Lorentz-boost proposition of the notes' Lorentzian section, which version~3 listed as an exception, is correct. Appendix~\ref{app:changes} lists every change.

\subsection*{Organization}

Section~\ref{sec:problem} states the problem and what is known. Section~\ref{sec:corpus} maps the workbench's documents, with the exact scope of the supplement's no-two-shears theorem and the quartic encoding of solutions. Sections~\ref{sec:core} and~\ref{sec:families} contain the core arithmetic and the records, families and census. Section~\ref{sec:items} treats the September continuation and the monodromy covers, and Section~\ref{sec:notes} the author's notes, including the mock theta results. Section~\ref{sec:negative} collects the proved negative results with their scope, Section~\ref{sec:overlaps} the bridges and the directions they open, and Section~\ref{sec:open} the questions and what was not checked. The appendices give the catalogue of statements, the verification, and the changes of version~4.
